\documentclass{amsart}
% For dealing with references we use the comment environment
\usepackage{verbatim}
\newenvironment{reference}{\comment}{\endcomment}
%\newenvironment{reference}{}{}
\newenvironment{slogan}{\comment}{\endcomment}
\newenvironment{history}{\comment}{\endcomment}

% For commutative diagrams we use Xy-pic
\usepackage[all]{xy}

% We use 2cell for 2-commutative diagrams.
\xyoption{2cell}
\UseAllTwocells

% We use multicol for the list of chapters between chapters
\usepackage{multicol}

% This is generally recommended for better output
\usepackage{lmodern}
\usepackage[T1]{fontenc}

% For cross-file-references

% Package for hypertext links:
\usepackage{hyperref}

% For any local file, say "hello.tex" you want to link to please

% Theorem environments.
%
\theoremstyle{plain}
\newtheorem{theorem}[subsection]{Theorem}
\newtheorem{proposition}[subsection]{Proposition}
\newtheorem{lemma}[subsection]{Lemma}

\theoremstyle{definition}
\newtheorem{definition}[subsection]{Definition}
\newtheorem{example}[subsection]{Example}
\newtheorem{exercise}[subsection]{Exercise}
\newtheorem{situation}[subsection]{Situation}

\theoremstyle{remark}
\newtheorem{remark}[subsection]{Remark}
\newtheorem{remarks}[subsection]{Remarks}

\numberwithin{equation}{subsection}

% Macros
%
\def\lim{\mathop{\mathrm{lim}}\nolimits}
\def\colim{\mathop{\mathrm{colim}}\nolimits}
\def\Spec{\mathop{\mathrm{Spec}}}
\def\Hom{\mathop{\mathrm{Hom}}\nolimits}
\def\Ext{\mathop{\mathrm{Ext}}\nolimits}
\def\SheafHom{\mathop{\mathcal{H}\!\mathit{om}}\nolimits}
\def\SheafExt{\mathop{\mathcal{E}\!\mathit{xt}}\nolimits}
\def\Sch{\mathit{Sch}}
\def\Mor{\mathop{\mathrm{Mor}}\nolimits}
\def\Ob{\mathop{\mathrm{Ob}}\nolimits}
\def\Sh{\mathop{\mathit{Sh}}\nolimits}
\def\NL{\mathop{N\!L}\nolimits}
\def\CH{\mathop{\mathrm{CH}}\nolimits}
\def\proetale{{pro\text{-}\acute{e}tale}}
\def\etale{{\acute{e}tale}}
\def\QCoh{\mathit{QCoh}}
\def\Ker{\mathop{\mathrm{Ker}}}
\def\Im{\mathop{\mathrm{Im}}}
\def\Coker{\mathop{\mathrm{Coker}}}
\def\Coim{\mathop{\mathrm{Coim}}}

% Boxtimes
%
\DeclareMathSymbol{\boxtimes}{\mathbin}{AMSa}{"02}

%
% Macros for moduli stacks/spaces
%
\def\QCohstack{\mathcal{QC}\!\mathit{oh}}
\def\Cohstack{\mathcal{C}\!\mathit{oh}}
\def\Spacesstack{\mathcal{S}\!\mathit{paces}}
\def\Quotfunctor{\mathrm{Quot}}
\def\Hilbfunctor{\mathrm{Hilb}}
\def\Curvesstack{\mathcal{C}\!\mathit{urves}}
\def\Polarizedstack{\mathcal{P}\!\mathit{olarized}}
\def\Complexesstack{\mathcal{C}\!\mathit{omplexes}}
% \Pic is the operator that assigns to X its picard group, usage \Pic(X)
% \Picardstack_{X/B} denotes the Picard stack of X over B
% \Picardfunctor_{X/B} denotes the Picard functor of X over B
\def\Pic{\mathop{\mathrm{Pic}}\nolimits}
\def\Picardstack{\mathcal{P}\!\mathit{ic}}
\def\Picardfunctor{\mathrm{Pic}}
\def\Deformationcategory{\mathcal{D}\!\mathit{ef}}

\setlength{\emergencystretch}{3em}
\begin{document}
\title[Stacks Project, Tag 0EY6]{323 --- Unique reflexive finite extension of etale covers across depth-two complements}
\maketitle
\noindent Copyright (C) 2005--2025 Johan de Jong. Stacks Project authors. Source: \href{https://stacks.math.columbia.edu/tag/0EY6}{Tag 0EY6}.

\noindent Editorial difficulty note (not author text). Difficulty H2: The proof must extend a finite etale algebra to a coherent reflexive algebra and prove uniqueness among finite depth-two extensions. It also identifies etaleness with flatness using codimension-one ramification control and establishes flat-base-change compatibility of the extension construction..

\noindent Editorial provenance note (not author text). History: excerpt from revision \texttt{a0444\allowbreak 6e57e\allowbreak c1fbc\allowbreak 252a8\allowbreak 71afc\allowbreak ec775\allowbreak 2fb28\allowbreak 07b14}, pione.tex, lines 5219--5301. Reference presentation was adapted; original-author context and core support blocks were retained or added. The mathematical wording is unchanged. Licensed under GNU FDL 1.2 or later; no invariant sections or cover texts. See ../licenses/COPYING. Context blocks reproduce author definitions and supporting results; they are not additional counted problems.

\begin{lemma}
\label{lemma-extend-S2}
Let $j : U \to X$ be an open immersion of locally Noetherian schemes
such that $\text{depth}(\mathcal{O}_{X, x}) \geq 2$ for $x \not \in U$.
Let $\pi : V \to U$ be finite \'etale. Then
\begin{enumerate}
\item $\mathcal{B} = j_*\pi_*\mathcal{O}_V$ is a reflexive coherent
$\mathcal{O}_X$-algebra, set $Y = \underline{\Spec}_X(\mathcal{B})$,
\item $Y \to X$ is the unique finite morphism such that
$V = Y \times_X U$ and $\text{depth}(\mathcal{O}_{Y, y}) \geq 2$
for $y \in Y \setminus V$,
\item $Y \to X$ is \'etale at $y$ if and only if $Y \to X$ is flat at $y$, and
\item $Y \to X$ is \'etale if and only if $\mathcal{B}$
is finite locally free as an $\mathcal{O}_X$-module.
\end{enumerate}
Moreover, (a) the construction of $\mathcal{B}$ and $Y \to X$ commutes
with base change by flat morphisms $X' \to X$ of locally Noetherian
schemes, and (b) if $V' \to U'$ is a finite \'etale morphism with
$U \subset U' \subset X$ open which restricts to $V \to U$ over $U$,
then there is a unique isomorphism $Y' \times_X U' = V'$ over $U'$.
\end{lemma}

\begin{proof}
Observe that $\pi_*\mathcal{O}_V$ is a finite locally free
$\mathcal{O}_U$-module, in particular reflexive.
By Divisors, Lemma \href{https://stacks.math.columbia.edu/tag/0EBJ}{lemma reflexive S2 extend (Tag 0EBJ)}
the module $j_*\pi_*\mathcal{O}_V$ is the unique
reflexive coherent module on $X$ restricting to
$\pi_*\mathcal{O}_V$ over $U$. This proves (1).

\medskip\noindent
By construction $Y \times_X U = V$.
Since $\mathcal{B}$ is coherent, we see that $Y \to X$ is finite.
We have $\text{depth}(\mathcal{B}_x) \geq 2$ for $x \in X \setminus U$
by Divisors, Lemma \href{https://stacks.math.columbia.edu/tag/0EBI}{lemma reflexive S2 (Tag 0EBI)}.
Hence $\text{depth}(\mathcal{O}_{Y, y}) \geq 2$ for $y \in Y \setminus V$
by Algebra, Lemma \href{https://stacks.math.columbia.edu/tag/0AUK}{lemma depth goes down finite (Tag 0AUK)}.
Conversely, suppose that $\pi' : Y' \to X$ is a finite morphism such that
$V = Y' \times_X U$ and $\text{depth}(\mathcal{O}_{Y', y'}) \geq 2$
for $y' \in Y' \setminus V$. Then $\pi'_*\mathcal{O}_{Y'}$
restricts to $\pi_*\mathcal{O}_V$ over $U$ and satisfies
$\text{depth}((\pi'_*\mathcal{O}_{Y'})_x) \geq 2$ for
$x \in X \setminus U$ by
Algebra, Lemma \href{https://stacks.math.columbia.edu/tag/0AUK}{lemma depth goes down finite (Tag 0AUK)}.
Then $\pi'_*\mathcal{O}_{Y'}$ is canonically isomorphic
to $j_*\pi_*\mathcal{O}_V$ for example by
Divisors, Lemma \href{https://stacks.math.columbia.edu/tag/0E9I}{lemma depth 2 hartog (Tag 0E9I)}.
This proves (2).

\medskip\noindent
If $Y \to X$ is \'etale at $y$, then $Y \to X$ is flat at $y$.
Conversely, suppose that $Y \to X$ is flat at $y$.
If $y \in V$, then $Y \to X$ is \'etale at $y$.
If $y \not \in V$, then we check (1), (2), (3), and (4) of
Lemma \href{https://stacks.math.columbia.edu/tag/0BJH}{lemma ramification quasi finite flat (Tag 0BJH)} hold
to see that $Y \to X$ is \'etale at $y$. Parts (1) and (2)
are clear and so is (3) since $\text{depth}(\mathcal{O}_{Y, y}) \geq 2$.
If $y' \leadsto y$ is a specialization and $\dim(\mathcal{O}_{Y, y'}) = 1$,
then $y' \in V$ since otherwise the depth of this local ring
would be $2$ a contradiction by
Algebra, Lemma \href{https://stacks.math.columbia.edu/tag/00LK}{lemma bound depth (Tag 00LK)}.
Hence $Y \to X$ is \'etale at $y'$ and we conclude (4) of
Lemma \href{https://stacks.math.columbia.edu/tag/0BJH}{lemma ramification quasi finite flat (Tag 0BJH)} holds too.
This finishes the proof of (3).

\medskip\noindent
Part (4) follows from (3) and the fact that $((Y \to X)_*\mathcal{O}_Y)_x$
is a flat $\mathcal{O}_{X, x}$-module if and only if $\mathcal{O}_{Y, y}$
is a flat $\mathcal{O}_{X, x}$-module for all $y \in Y$ mapping to $x$, see
Algebra, Lemma \href{https://stacks.math.columbia.edu/tag/00HT}{lemma flat localization (Tag 00HT)}. Here we also
use that a finite flat module over a Noetherian ring is finite locally
free, see Algebra, Lemma \href{https://stacks.math.columbia.edu/tag/00NX}{lemma finite projective (Tag 00NX)}
(and
Algebra, Lemma
\href{https://stacks.math.columbia.edu/tag/00FP}{lemma Noetherian finite type is finite presentation (Tag 00FP)}).

\medskip\noindent
As to the final assertions of the lemma, part (a) follows from
flat base change, see
Cohomology of Schemes, Lemma \href{https://stacks.math.columbia.edu/tag/02KH}{lemma flat base change cohomology (Tag 02KH)}
and part (b) follows from the uniqueness in (2) applied to the restriction
$Y \times_X U'$.
\end{proof}
\end{document}
